\section{Discussion}\label{sec:discussion}
The observation process has two roles in this problem. Nonresolution creates an unidentified terminal harmful mass, while the resolved path supplies an observable information requirement. The profile decomposition separates these roles: $\tau_\Pi$ accounts for dependence already observed in resolved timing, leaving a smaller allowance for terminal selection. Coarsening removes part of that observable requirement, leaving more dependence available to adverse terminal completions. This accounts for the larger endpoint risk set under a common budget.

The exact frontier $\rho_\Pi^*(\beta)$ provides a direct sensitivity report. It identifies the budgets at which a particular risk decision follows from the observed law, together with the budgets ruled out by compatibility. Its shift under refinement is the same conditional information that appears in the profile identity. The population decision example and the endpoint-uncertifiable power scenario give two realizations of this geometry, at different risk and information budgets.

Finite-sample resolution is a design choice with a cost. The strongest timing law supports markedly earlier certification with two bands, whereas its matched null favors endpoint-only inference. Additional timing bands can tighten the population bound while increasing confidence-region uncertainty and numerical relaxation. These findings support choosing a prespecified resolution appropriate to the expected timing structure; they do not support selecting the most favorable certificate after inspecting multiple unadjusted analyses.

The present guarantees concern trustworthy labels and event records, a fixed terminal horizon, a stable categorical law, and a sensitivity assumption that holds for the target stream. The experiments demonstrate the mechanism on designed synthetic laws. They do not estimate its prevalence across applications or establish a portable calibration of $\rho$. Richer joint confidence regions, tighter projection envelopes and application-specific sensitivity elicitation are natural questions, but changes to these components would require their own validation. The current computational contribution is a transparent conservative implementation with traceable numerical evidence, rather than a finite-sample optimality claim.

\section{Conclusion}\label{sec:conclusion}
Under a mutual-information constraint on a categorical outcome-resolution path, latent harmful risk has a sharp interval characterized by one scalar completion. Observable timing information determines the compatibility floor and shifts the risk-certification frontier exactly under refinement. Projecting a simultaneous observable-law confidence sequence through this model gives an anytime-valid certificate with conservative numerical computation. The verified experiments show both consequences of retaining timing: decisions unavailable to an endpoint population analysis, and high finite-sample certification at substantially shorter horizons under informative timing. The matched control identifies the condition behind that benefit: timing must add enough relevant information to offset the statistical cost of observing it at finer resolution.
